\section{Method}
\label{sec:method}

\subsection{Geometry encoding and point parameterization}

The encoder receives point coordinates together with first and second
derivatives normalized by chord-length parameter gaps. It produces local tokens $F\in\R^{M\times h}$ and a pooled curve
descriptor $g\in\R^h$. A parameter head predicts positive intervals, which are
normalized to form a strictly increasing parameter sequence. A positional
encoding of the predicted parameters is added to the local tokens before knot
decoding.

\subsection{Canonical structural supervision}

Starting from a generated source knot vector, we greedily remove the single
knot whose removal gives the smallest refitting error. A removal is accepted
when the normalized Euclidean RMS remains below a prescribed tolerance
$\epsilon$. The terminal representation defines the canonical count $K^*$ and
knot locations $U^*$. Thus the target describes required geometric complexity
rather than the arbitrary representation used to generate the samples.

\subsection{Ordinal knot-count prediction}

A set of count queries cross-attends to the position-encoded local tokens and
produces a scalar complexity evidence $e$. Ordered thresholds $b_r$ define
\begin{equation}
  s_r=P(K\ge r)=\sigma(e-b_r),\qquad r=1,\ldots,\Kmax.
\end{equation}
Adjacent survival probabilities define a normalized class distribution. The
head is trained using binary ordinal targets $\mathbf 1[K^*\ge r]$ and an
asymmetric expected-overcount penalty.

\subsection{Count-conditioned knot decoder}

All admissible counts share interval queries, cross-attention, and output
weights. A learned count embedding conditions the first $K+1$ interval queries.
Positive interval logits are normalized to obtain exactly $K$ strictly ordered
internal knots. During training, the representation indexed by the true count
receives the knot-position and fitting losses.
The direct neural prediction used for validation is
\begin{equation}
  \widehat K=\arg\max_K p_K,
  \qquad \widehat U=\mathcal D_{\widehat K}(F,g).
\end{equation}
The deployment refinement described below may instead select among complete
count branches. No candidate deletion, binary gate, or probability threshold
is used.

\subsection{Training objective}

For a sample with count $K^*$ and true internal knots $U^*$, the objective is
\begin{equation}
  \mathcal L=\lambda_f\mathcal L_{fit}
  +\lambda_t\mathcal L_t
  +\lambda_K\mathcal L_{count}
  +\lambda_o\mathcal L_{over}
  +\lambda_u\mathcal L_{knot},
\end{equation}
where
\begin{align}
  \mathcal L_{fit}&=M^{-1}\sum_i\|C(\widehat t_i)-q_i\|_2^2,\\
  \mathcal L_t&=M^{-1}\sum_i(\widehat t_i-t_i^*)^2,\\
  \mathcal L_{count}&=\Kmax^{-1}\sum_r
  \operatorname{BCEWithLogits}(e-b_r,\mathbf 1[K^*\ge r]),\\
  \mathcal L_{over}&=\max(0,\mathbb E[K]-K^*),\\
  \mathcal L_{knot}&=(K^*)^{-1}\sum_{j=1}^{K^*}
  \operatorname{SmoothL1}(\widehat u_j^{(K^*)}-u_j^*).
\end{align}

\subsection{Standard B-spline deployment}

Every complete count-conditioned representation defines an open knot vector.
Equation \eqref{eq:control_solve} is solved for each admissible count, after
which BIC augmented by the learned count prior selects the deployed model.
This compares complete spline orders and does not delete individual candidates.

\begin{figure}[t]
  \centering
  \fbox{\parbox[c][35mm][c]{0.90\linewidth}{\centering
  \TODO{Architecture figure: ordered samples $\rightarrow$ encoder and
  parameter head $\rightarrow$ CountHead $\rightarrow$ selected
  count-conditioned decoder $\rightarrow$ standard B-spline refit.}}}
  \caption{Overview of the proposed count-conditioned fitting pipeline.}
  \label{fig:architecture}
\end{figure}
